\section{MA Thesis Guidelines}

\subsection{Typesetting Tool}
In general, we highly recommend the use of \LaTeX{} for almost all MA theses as this will take care of almost all formatting and style aspects mentioned below. A tutorial on learning \LaTeX{} can be found here: \href{https://www.overleaf.com/learn/latex/Learn_LaTeX_in_30_minutes}{Link}.

\subsection{Length \& Structure}

The scope of the MA thesis should be defined by the topic. Typically, the number of pages should range anywhere between 10-20 written pages net, gross of further pages such as the title page, table of contents, references, table of figures, table of tables and further appendices.

Your thesis should have a logical and self-explanatory structure. Typically, your sections could be structured as follows:
\begin{itemize}
\item \textbf{Title Page}: Title, possibly subtitle, author, supervisor, school, class name, ``Maturitätsarbeit'', place and date
\item \textbf{Table of Contents}: Overview of the sections and subsections of your thesis (in \LaTeX: just one line \lstinline|\tableofcontents|).
\item \textbf{Introduction} (numbering applies from here: 1(.1, .2), 2(.2,.3) etc.): Question / objective of your thesis, motivations, approach.
\item \textbf{Methodology}: Overview of possible approaches to address your question and in-depth presentation of your approach, \textit{without}(!) presenting or discussing any results. If possible, also define a success metric (such as a number which you are trying to maximize, but can also be subjective) to assess the success of your thesis.
\item \textbf{Results}: present the results of your thesis
\item \textbf{Discussion}: Assess and discuss the results. 
\item \textbf{Outlook}: What could or should be done to further improve your work if more time was available? Don't try to be too self-critical here but rather, this should reflect the fact that there is always more work to do (and in the case of scientific works, grants to obtain ;))
\item \textbf{Conclusion}: A summary of the most important results and personal experiences during the process. Any further thoughts regarding the topic can also be inserted here.
\item \textbf{Appendices} (numbered Appendix A(.1, .2), B(.1,.2), etc.): Any appendix documents, such as questionnaires, code listings, further figures or tables, documents etc.
\item \textbf{References}: Listing of all Materials, sources including online sources (in \LaTeX: just one line \lstinline|\printbibliography|, have a look at the package \lstinline|biblatex|: \href{https://de.overleaf.com/learn/latex/Bibliography_management_with_biblatex}{Link})
\item \textbf{Table of Figures} (in \LaTeX: just one line \lstinline|\tableoffigures|)
\item \textbf{Table of Table} (in \LaTeX: just one line \lstinline|\tableoftables|)
\item \textbf{Glossary}: List of abbreviations, if many abbreviations are used in the text  (in \LaTeX: just one line \lstinline|\printglossaries|, have a look at the package \lstinline|glossaries|: \href{https://www.overleaf.com/learn/latex/Glossaries}{Link})

\end{itemize}

\subsection{Formal Aspects \& Formatting}

If you'd still like to use Word or another typesetting tool, please stick to the following general rules:



\begin{itemize}

\item Justification: Block justification with hyphenation
\item Page margins: Roughly 3 cm left and right
\end{itemize}

\subsection{Quotation Style}
You should clarify all contributions to your topic by citing any sources appropriately in your text. Again, citations are greatly facilitated if you use \LaTeX. We recommend reading through the following tutorial in depth: \href{https://de.overleaf.com/learn/latex/Bibliography_management_with_biblatex}{Link}. When you quote something in the middle of a sentence, put the citation right after the quote. In all other cases, meaning sentences which are based on a source but do not contain literal quotes, add the citation in the end.

Please use the following parameters for your citations: 

\begin{lstlisting}
\usepackage[maxbibnames=99,sorting=none,style=numeric,backend=biber]{biblatex}
\end{lstlisting}


\subsection{Writing style}

Prefer passive sentences and avoid the words ``I'' and ``we''. This is standard in scientific writing, as it forces you into a more objective way of thinking.

\begin{myexample}
Instead of writing:
\begin{quotation}
After I searched the internet, I found that there were many options for doing \textit{xyz}...
\end{quotation}
You could write:
\begin{quotation}
In order to achieve \textit{xyz}, it seems that a variety of approaches are possible (insert citations!)
\end{quotation}
\end{myexample}

In order to check your grammar, there are lots of good tools and approaches:
\begin{itemize}
\item Have your MA thesis proofread by at least 2-3 close acquaintences or family members who are proficient in English.
\item Use online tools such as \url{https://www.grammarly.com/grammar-check} or ChatGPT (but don't write your entire thesis using the tool – this greatly hampers creativity and there can and will be checks for plagiarism).
\end{itemize}

Feel free to suggest any clarifications and additions to these guidelines by contacting \href{\mainauthormail}{\mainauthor}.